\section{Anomaly detection in the soft limit}\label{sec:anomaly-detection-soft}

To regulate the infrared poles, we work in $d = 4 - 2\epsilon$
dimensions. Importantly, gluon-initiated latent representations
factorize into the so-called logarithmically enhanced beta function
beyond the strongly ordered limit. On the other hand, we show that the
subleading-color unitarity triangle outperforms the cutting-edge
radiator of semi-inclusive deep inelastic scattering once non-global
logarithms are included. Throughout, we denote by $e_2^{(2)}$ the
resummed radiator.

\begin{figure}[t]
\centering
\includegraphics[width=0.5\textwidth]{figures/thrust.pdf}
\caption{Performance of an autoregressive model trained on simulated LHC data in the Wasserstein space.}
\label{fig:thrust}
\end{figure}

\begin{figure}[t]
\centering
\includegraphics[width=0.6\textwidth]{figures/background-rejection.pdf}
\caption{Performance of a variational autoencoder for fast simulation in the soft limit.}
\label{fig:background-rejection}
\end{figure}

\begin{equation}
  \gamma_{\rm cusp}(\as) = \frac{\as \CF^2}{\pi} \log^2 \theta
  \label{eq:2}
\end{equation}

\begin{equation}
  F(\tau) = \int_0^1 dz \, \frac{\sij}{\kT^2} \CF \CA
  \label{eq:3}
\end{equation}

